\subsection{Três entradas para o mesmo pagamento}
\label{sec:tres-entradas-equilibrio}
\index{pseudo-compromisso@pseudo-compromisso!escolha das mascaras@escolha das máscaras}

No princípio Alice tinha as onze unidades numa só saída. Imaginemos
que as recebeu em três ocasiões diferentes, 2 da primeira vez, 4 da
segunda e 5 da terceira. Bob há-de receber os mesmos 7, o troco
continua a ser 3 e a taxa 1. O que muda na conta do nó? Há agora
três entradas, cada qual com o seu anel e a sua imagem de chave.
No entanto, a soma dos montantes de entrada só tem de satisfazer
uma igualdade. Esta diferença entre três autorizações e uma única
conta explica a liberdade que a carteira tem ao escolher as máscaras.

Para acompanhar a álgebra até ao fim, demos valores pequenos às
máscaras do exemplo. Os geradores são \(G,H\) no subgrupo de ordem
\(l\). Os números que seguem servem apenas para expor as operações.
Alice recebeu as três saídas com os compromissos
\[
       C_1=4G+2H,\qquad C_2=9G+4H,\qquad C_3=12G+5H .
\]
Nas saídas novas suponhamos que as máscaras derivadas para Bob e
para o troco valem 6 e 8. Do lado que paga, o nó terá então de fazer
a soma
\[
               (6G+7H)+(8G+3H)+H=14G+11H .
\]
Se somássemos os compromissos recebidos por Alice, obteríamos
\(25G+11H\). Ambas as expressões representam as onze unidades, e
no entanto há uma diferença de \(11G\). As máscaras nas saídas
antigas nasceram em pagamentos independentes, cujos autores
desconheciam o que Alice havia de fazer hoje.

Chamemos \(u_1,u_2,u_3\) às três máscaras dos pseudo-compromissos.
Para comparar as somas basta que satisfaçam, em \(\mathbb F_l\),
\[
                              u_1+u_2+u_3=14 .
\]
Por exemplo, escolhendo \(u_1=3\) e \(u_2=5\), resta
\(u_3=14-3-5=6\). Ficam \(\widetilde C_1=3G+2H\),
\(\widetilde C_2=5G+4H\) e \(\widetilde C_3=6G+5H\).
As três autorizações usam, nas respectivas posições secretas, as
diferenças que se podem ler na última coluna.
\[
\begin{array}{c|c|c|c}
 j & C_j &\widetilde C_j&C_j-\widetilde C_j \\ \hline
 1&4G+2H &3G+2H&G\\
 2&9G+4H &5G+4H&4G\\
 3&12G+5H&6G+5H&6G
\end{array}
\]
Na primeira entrada a testemunha auxiliar de CLSAG vale 1, na
segunda vale 4 e na terceira vale 6. Cada uma acompanha a chave
privada de gasto da sua própria saída. Vejamos o que o nó consegue
refazer desta tabela. Para cada um dos dezasseis membros de um anel
conhece um compromisso histórico e pode subtrair-lhe o
pseudo-compromisso dessa entrada. A assinatura atesta o conhecimento
das duas testemunhas na mesma posição. A coluna que aqui exibe os
três resultados reais esteve à nossa disposição porque acompanhámos
as escolhas da carteira. Ao nó chegam três listas de dezasseis
diferenças, juntamente com a autorização de cada uma.

Para conferir o balanço, o nó já dispõe dos três
pseudo-compromissos sem procurar membro algum dos anéis. A soma é
\[
 (3G+2H)+(5G+4H)+(6G+5H)
      =(3+5+6)G+(2+4+5)H=14G+11H .
\]
É o ponto que obteve do lado das novas saídas e da taxa. As imagens
de chave têm outra função, que se conserva mesmo se duas listas de
dezasseis membros partilharem alguns candidatos. O nó verifica as
três imagens e o seu estado de utilização. Uma coincidência entre
candidatos de dois anéis é possível. Cada entrada tem a sua imagem
e o seu fecho.

Vale a pena contar o que a igualdade das máscaras determina.
Tendo \((3,5,6)\) por solução, todo o triplo
\[
                  (3+s,\ 5+t,\ 6-s-t),\qquad
                  s,t\in\mathbb F_l ,
\]
dá a mesma soma. Reciprocamente, quaisquer máscaras com soma 14
determinam \(s=u_1-3\), \(t=u_2-5\), e a última componente fica
fixada. Eis uma correspondência entre os pares de escalares e as
soluções da equação. A equação linear, considerada em todo o corpo,
tem por isso \(l^2\) soluções. Uma regra concreta de amostragem
pode restringir os valores sorteados. Para conhecer a distribuição
com que a carteira percorre as soluções é preciso conhecer também
a lei desses sorteios.

A liberdade tem uma interpretação pública. Se passarmos de
\((3,5,6)\) para \((5,5,4)\), o primeiro pseudo-compromisso recebe
\(2G\), o último perde \(2G\), e a soma fica intacta. Em contrapartida,
as listas de diferenças dos anéis correspondentes mudam, e cada
assinatura tem de ser construída para o pseudo-compromisso efectivamente
declarado. A igualdade dos totais permite várias escolhas antes da
assinatura. A mensagem e o transcrito fixam as escolhas efectuadas
para aquela transacção.

Com \(m\) entradas, a mesma conta deixa \(m-1\) coordenadas livres.
Fixada a soma \(Y\) das máscaras das saídas, a última coordenada é
\[
                 u_m=Y-\sum_{j=1}^{m-1}u_j .
\]
A uma entrada, a soma vazia vale zero, o que explica por que razão a
carteira de Alice, no primeiro exemplo, tinha uma única máscara possível para
o pseudo-compromisso. A três entradas tem duas escolhas preliminares.
Em ambos os casos a última vem da mesma igualdade. Cada entrada
precisa da sua autorização. Passar de uma para três pede
mais duas assinaturas e mais duas imagens, ainda que os montantes
entregues a Bob e a Alice fiquem exactamente nos mesmos compromissos.
A prova de intervalo abrange as duas saídas novas.
